\section{Justificativa da precificação}
\label{sec:justificativa}

O preço foi definido a partir dos três critérios da atividade: margem de lucro desejada, valor
percebido pelo cliente e comparação com a concorrência.

\subsection{Margem de lucro desejada}

A meta do grupo, ainda uma hipótese a validar, é uma \textbf{margem de contribuição consolidada de pelo menos 50\%} e
uma \textbf{margem líquida de cerca de 15\%} no cenário enxuto, no mês 12. O modelo estimado
atinge 50,5\% e 14,8\% (Tabela~\ref{tab:consolidado}), respectivamente, com contribuição de
R\$~2.232,15 e resultado de R\$~756,32, incluindo os R\$~1.500,00 de convênio (estimativa do grupo).

\begin{center}
\begin{minipage}{0.92\textwidth}
\textbf{Quadro auxiliar — meta de contribuição por unidade versus modelo estimado.}
Metas definidas como hipóteses do grupo.\par\medskip
\begin{tabular}{L{6.0cm} L{3.0cm} R{3.5cm}}
\toprule
\cabfundo \cab{Unidade} & \cab{Meta estimada} & \cab{Modelo} \\
\midrule
Assinatura mensal & $\geq70\%$ & 76,7\% (R\$~9,89) \\
Plano Parceiro mensal & $\geq80\%$ & 89,0\% (R\$~71,13) \\
Atendimento antes do fundo & $\geq60\%$ & 68,1\% (R\$~5,24) \\
Atendimento após o fundo & $\geq0\%$ & $-17{,}6\%$ (R\$~$-1{,}36$) \\
\bottomrule
\end{tabular}
\end{minipage}
\end{center}

As metas são hipóteses, não resultados validados. São margens de contribuição, e não lucro
líquido por unidade, pois não incluem rateio dos custos fixos. No atendimento, os 68,1\%
valem apenas \textbf{antes do fundo}. Após os R\$~6,60 de liquidação, a contribuição é
negativa; a meta de não perder por atendimento não é atendida. No mês 12, o resultado
consolidado só é positivo com convênio/patrocínio, juntamente com as demais receitas.

O preço definido de R\$ 12,90 garante uma margem de contribuição unitária de 76,7\% (R\$ 9,89 por assinante), cobrindo as taxas de loja, impostos e a alocação de suporte por usuário.

O fundo de R\$~990,00 \textbf{já foi descontado} da contribuição consolidada de R\$~2.232,15.
Sem convênios, a contribuição cai para R\$~822,15 e o resultado A para R\$~$-653{,}68$;
o equilíbrio sobe de 303 para 2.876 usuários (Tabela~\ref{tab:equilibrio}).
Se, em vez da premissa de 6\%, for aplicada a alíquota nominal inicial de 15,5\% do
Anexo V, a estimativa do resultado A cai para R\$~336,28 (margem líquida de cerca de 7,6\%),
como detalhado na análise de sensibilidade tributária da Seção~\ref{sec:consolidado}.

\subsection{Valor percebido pelo cliente}

\begin{itemize}
  \item \textbf{Para o tutor assinante:} R\$~12,90/mês somam R\$~154,80/ano, valor próximo ao
  extremo inferior de uma consulta em Belo Horizonte (R\$~150 a R\$~350). Em troca, recebe
  lembretes de vacinas, múltiplos pets, prontuário ampliado e descontos nos parceiros.
  \item \textbf{Para o tutor de baixa renda:} o núcleo é gratuito. O benefício percebido é pagar
  menos pelo atendimento (preço social e créditos) e evitar emergências por falta de prevenção,
  a dor validada nas entrevistas.
  \item \textbf{Para a clínica:} R\$~79,90/mês devem ser compensados por novos clientes e pelo
  preenchimento de horários ociosos. A hipótese do grupo, ainda a validar, de que um atendimento adicional por
  semana cubra a mensalidade depende da \textbf{contribuição líquida da clínica}, após seus
  custos variáveis, impostos e eventuais taxas, e não do ticket bruto de R\$~110,00.
  A contribuição incremental mensal precisa ser de pelo menos R\$~79,90; sua confirmação
  exige entrevistas e dados de custos dos parceiros, ainda a confirmar.
  \item \textbf{Limite:} a disposição a pagar das classes C, D e E ainda não foi medida. O teste de
  aceitação (Capítulo~\ref{sec:teste}) existe para validar esse ponto.
\end{itemize}

\subsection{Análise da concorrência}

\begin{longtable}{L{2.9cm} L{3.1cm} L{2.8cm} L{4.2cm}}
\caption{Concorrência e alternativas: preços indicativos.}\label{tab:concorrencia}\\
\toprule
\cabfundo \cab{Alternativa} & \cab{Oferta} & \cab{Preço} & \cab{Observação} \\
\midrule
\endfirsthead
\toprule
\cabfundo \cab{Alternativa} & \cab{Oferta} & \cab{Preço} & \cab{Observação} \\
\midrule
\endhead
\bottomrule
\endfoot
Petlove Saúde, plano Leve & Consultas e vacinas essenciais & R\$~14,90 a R\$~17,90/mês & Preço indicativo, consultado em 05/10/2026 em site intermediário; a confirmar e sujeito a variação \\
Petlove Saúde, plano Tranquilo & Vacinas, clínico geral e urgência & R\$~34,90 a R\$~39,90/mês & Preço indicativo, consultado em 05/10/2026 em site intermediário; a confirmar e sujeito a variação \\
Petlove Saúde, planos Ideal e Essencial & Especialistas, cirurgias e internação & R\$~89,90 e R\$~194,90/mês & Preço indicativo, consultado em 05/10/2026 em site intermediário; a confirmar e sujeito a variação \\
Planos de saúde pet em geral & Cobertura por faixa & Entrada: R\$~20 a R\$~70/mês & Faixa de mercado em 2026 \\
Flockr Premium & Aplicativo de registro e lembretes do pet & R\$~9,90 a R\$~12,90/mês & Preço indicativo, consultado em 05/10/2026 em site intermediário; a confirmar e sujeito a variação \\
Atendimento particular em BH & Consulta clínica geral & R\$~150 a R\$~350 & Sem desconto social \\
Vacinas no particular & V10 e antirrábica & V10: R\$~80 a R\$~150; antirrábica: R\$~60 a R\$~80 & Antirrábica é gratuita em campanhas públicas \\
\end{longtable}

\textbf{Posicionamento.} O PetCuida+ a R\$~12,90 fica no mesmo patamar de aplicativos de registro do
pet e abaixo do plano Tranquilo da Petlove \textit{nas faixas indicativas da
Tabela~\ref{tab:concorrencia}, ainda a confirmar}. O \projeto{} \textbf{não é um plano de saúde}: não
oferece cobertura nem rede de seguro, e sim prevenção, triagem orientativa, créditos solidários e
tabela social. Por isso compete pelo orçamento do tutor de baixa renda, que dificilmente contrata
planos de mensalidade mais alta. Os valores de planos foram consultados em 05/10/2026 em
sites intermediários, sujeitos a variação por região e plano, conforme o levantamento do grupo.
Nesta revisão não foi possível confirmar os preços da Petlove e do Flockr em fontes oficiais
acessíveis. Por isso, as faixas não são apresentadas como cotações oficiais verificadas.